\documentclass[12pt]{article}

\input{.house-style/preamble.tex}

\hypersetup{
  pdftitle={Bridge Principles without UG: Projectibility, Crosslinguistic Evidence, and Linguistic Ontology},
  pdfkeywords={Universal Grammar, projectibility, linguistic ontology, crosslinguistic evidence, bridge principles, comparative method}
}

\setlength{\emergencystretch}{2em}

\title{Bridge Principles without UG: Projectibility, Crosslinguistic Evidence, and Linguistic Ontology}
\author{Brett Reynolds \orcidlink{0000-0003-0073-7195}%
\thanks{Contact: \href{mailto:brett.reynolds@humber.ca}{brett.reynolds@humber.ca}}\\
Humber Polytechnic \& University of Toronto}
\date{\today}

\begin{document}
\maketitle

\begin{abstract}
Crosslinguistic argument isn't methodologically innocent. If evidence from one language is used to analyse another, some bridge principle is needed to explain why the evidence travels. This paper accepts that constraint but rejects the stronger claim that only a Chomskyan Universal Grammar can supply the bridge. Universal Grammar is one proposed licence for evidential travel, not the condition of possibility for linguistic comparison as such. A projectibility-based account instead asks what a proposed category lets us infer, what property-and-relation pattern supports that inference, and where the projection should be localized, demoted, or rejected. Reiss's own French-English example involving \mention{any} is reconstructed to distinguish a diagnostic cue, a scoped bridge, and a robust bridge; on its own, the French evidence reaches only cue level. Historical-comparative reconstruction then supplies a positive non-UG bridge: common descent and regular correspondence can license evidential transfer without appeal to UG. The result is a pluralist realism about linguistic targets: mental grammars, public standards, usage profiles, typological comparanda, and norm-governed varieties can all be real enough for inquiry when their projection targets and failure conditions are explicit.
\end{abstract}

\noindent\textbf{Keywords:} Universal Grammar; projectibility; linguistic ontology; crosslinguistic evidence; bridge principles; comparative method

\section{Introduction: Bridge Principles, Not Boundary Policing}

Comparing English, French, Japanese, Yoruba, or Warlpiri isn't the problem. The issue is what licenses a particular comparison. The starting concession is substantial: linguistic categories don't announce themselves from raw speech, text, or judgments, and evidence from one language doesn't automatically bear on another. Reiss is right about that burden \citep{reiss2026armchair}. The disagreement is with the monopoly claim that crosslinguistic evidence is licensed only by a genetically specified Universal Grammar.

Reiss states the strongest version of the claim in exactly these terms: the UG postulate \enquote{allows data from each language to bear on the analysis of all languages} \citep[sec.~1]{reiss2026armchair}, and \enquote{Anyone who denies UG has no grounds for importing analytic categories} \citep[sec.~6]{reiss2026armchair}. These aren't local claims about one successful research programme. They're exclusionary claims about the field's evidential rights.

The exclusionary claim has two faces. One is evidential: why should French evidence bear on English, or Japanese evidence on Icelandic? The other is classificatory: why are French, Warlpiri, and Swahili in the same domain of inquiry while chess, Python, and tango aren't? Reiss is right that neither question can be waved away. The mistake is to treat the need for a domain and the need for transfer as if both had already been answered by the word \emph{UG}.

Two failures undermine that monopoly claim. First, a general postulate of shared human linguistic endowment doesn't yet specify which inferences are licensed. It doesn't identify a projection target, a transfer condition, a scope limit, or a failure mode. Second, Reiss's argument moves from a programme-internal definition of \enquote{linguists} to a conclusion about what \term{linguists} as such are entitled to do. That move requires the very sort of bridge principle the paper says only UG can provide.

Projectibility supplies the alternative: what observing some features licenses us to expect about others \citep{Goodman1955,reynolds2026kindsProjectibilityProfiles}. A projectibility-first linguistics asks what a proposed linguistic kind lets us infer, which profile supports that inference, which diagnostics locate the profile, and what defeats the projection. Universal Grammar can be part of a proposed support story. It isn't the condition of possibility for linguistic comparison.

\section{What Reiss Gets Right}

Reiss's best point is methodological. Theoretical linguistics can be empirical without being experimental, corpus-based, or statistical. Good armchair work isn't mere introspection detached from the world. It's a form of theory-mediated empirical reasoning in which judgements, contrasts, paradigms, and crosslinguistic patterns are made to bear on abstract structure.

A second point follows. Categories such as word, sentence, negative-polarity item, segment, syllable, and interrogative element aren't read directly from the acoustic signal, the printed page, or a spreadsheet of corpus counts. A statistician still has to know what's being counted. A corpus linguist still has to decide what counts as a token of the relevant sort. A psycholinguistic experiment still has to decide which contrast the stimuli instantiate.

The third concession matters most. Crosslinguistic evidence needs a warrant. If French data is used to analyse English, or Japanese data is used to analyse English and Icelandic, then the inference isn't justified by the mere fact that all are called languages. Reiss's examples of empirical argumentation devices put pressure on any view that treats crosslinguistic comparison as automatic.

Nor should the reply pretend that UG-style research never provides content. A Chomskyan account can offer precise hypotheses about locality, binding, polarity, category features, movement, or possible dependencies, and those hypotheses can generate crosslinguistic expectations. The objection is narrower and stronger: successful UG-shaped bridges for I-language inquiry don't establish that UG is the only possible bridge, nor that every legitimate linguistic target has to be an I-language target.

Those concessions make the disagreement narrower. The reply accepts theory, realism, armchair method, and comparison. The objection is that Reiss identifies a real bridge-principle burden and then overstates what follows from it.

\section{The Monopoly Claim}

Vulnerability enters at the move from bridge principles to Universal Grammar as the only possible bridge. Reiss's chapter moves among several claims: languages are comparable; individual mental grammars are real objects; relevant primitives are innate and genetically specified; data from one language can bear on another; only Universal Grammar licenses that evidential travel. The last claim doesn't follow from the earlier ones.

That flaw is fatal to the monopoly claim, not to Chomskyan inquiry as such. A general postulate of shared human linguistic endowment may define a research domain for \enquote{linguists}, but it doesn't by itself say which comparison is licensed, what the projection target is, how far the evidence travels, or when the inference fails. A bridge that never specifies its failure conditions is too permissive to do the work Reiss assigns to it.

A familiar retreat doesn't help. If UG merely makes crosslinguistic evidence potentially relevant, the monopoly conclusion doesn't follow: potential relevance isn't a licence for any particular transfer. If, instead, a substantive UG theory licenses a particular transfer, then UG is no longer functioning as a bare domain-defining postulate. It's a first-order bridge with target, diagnostics, scope, and defeaters, and it has to be audited alongside other bridge principles. A domain postulate may mark a field; it doesn't carry evidence across every point within it. Reiss asks UG to do both jobs at once. A thin UG postulate can't license the monopoly; a thick UG theory is one auditable bridge among others.

Reiss also shifts roles for UG. It appears as an innate faculty, a universal inventory of primitives, a theory of general linguistics, a methodological postulate, a null hypothesis, and a licence for transfer. Those roles can't be interchanged without cost. A postulate may orient a research programme, but it doesn't by itself supply the primitives, diagnostics, or defeaters needed for a particular transfer. A substantive theory may supply those things, but then it's no longer a neutral precondition of inquiry. It's one bridge under audit.

The Newton analogy shows the problem. Newton's postulate grouped apples, tides, and planets by generating equations, measures, counterfactual expectations, and conditions under which the theory would be strained. Reiss's UG postulate is treated as if it generated jurisdiction: a password that admits preferred transfers and excludes rival licences. It's asked to do the grouping work, the transfer work, and the exclusionary work at once.

A substantive UG theory may do some of that work for a particular comparison. It may predict that a dependency should be constrained by locality, that an apparent lexical item should decompose into two grammatical objects, or that a surface contrast should reflect an unpronounced structural contrast. But then the warrant comes from the specified theory, diagnostics, and failure conditions, not from the bare fact that UG has been postulated. The examples show that Chomskyan \enquote{linguists} have found fruitful comparisons. They don't show that all legitimate comparisons require UG.

The empirical argumentation devices don't need to be rejected. Homophony, hidden structure, distributional contrast, and crosslinguistic prompts are ordinary forms of linguistic reasoning. Their legitimacy depends on the target and on the bridge that connects the prompt to the proposed analysis. A Chomskyan theory may provide such a bridge, but the device itself doesn't carry a UG credential. If the diagnostics, scope, and defeaters are supplied by functional recurrence, diachrony, contact history, or language-internal distribution, the argument remains empirical without becoming Reiss-style UG.

Nor does the domain question force the monopoly conclusion. Human languages can be grouped by a projectible profile rather than by a single genetic essence: productive pairings of form and interpretation; developmental uptake as a primary communicative system; externalization in speech, sign, writing, or other modalities; social normativity; diachronic transmission; variation; and recurrent possibilities of translation, paraphrase, error, correction, and grammatical judgement. The cluster is stabilized by biological endowment, acquisition and uptake, communicative function, social transmission, norm-governed reproduction, and repair under error.

That profile doesn't smuggle in the Chomskyan conclusion. Form-interpretation pairings and grammatical judgements presuppose grammaticality as a projectible property, not UG as its explanation. Even if a UG theory best explained part of the cluster's cohesion, that would help individuate the domain; it wouldn't license transfer from one language to another without a first-order bridge.

Chess, Python, and tango share some traits with human languages. They have conventions, learnability, sequencing, and in some cases formal syntax. But they don't support the same projection package: native-language acquisition, ordinary open-ended communicative uptake, historical sound and morphosyntactic change, indexical register effects, crosslinguistic translation, and grammar-internal acceptability contrasts don't travel to them as a set. They may be valuable comparanda. They aren't thereby co-domain targets for the same linguistic projections.

Projectibility-first linguistics can ask the missing questions directly. It can distinguish a claim about mental grammars from a claim about a public standard, a corpus profile, a typological comparandum, a diachronic pathway, or an institutional norm. Reiss's \enquote{linguists} can handle the first target on their own terms. They can't infer from that success that the other targets are illegitimate, or that \term{linguists} working on them lack grounds for comparison.

\section{Boundary Stipulation and Disciplinary Ontology}

Reiss's restriction is best granted rather than resisted. In what follows, I use \enquote{linguists}, with quotation marks, for the group his chapter stipulates: Chomsky and those who broadly adopt Chomsky's philosophical stance on the object and methods of inquiry. That local definition is legitimate. It names a research programme, not the field.

Presentation order matters. Before UG is credited with licensing field-wide evidential rights, the chapter has already marked some opponents as scholars who \enquote{self-identify as linguists} \citep[sec.~2]{reiss2026armchair}. The phrase has consequences: it turns colleagues into claimants and disagreement into a question of standing. That phrasing isn't the argument, but it marks the work the argument still owes. A boundary first drawn rhetorically still needs a bridge if it's to support a field-wide conclusion.

By \term{linguists}, in contrast, I mean researchers whose work is organized around systematic inquiry into human language phenomena. The category is projectible when it tracks a worldly pattern that licenses expectations about evidence, methods, training, questions, and products. The projected profile is sustained engagement with language data; explicit methods for identifying, comparing, and explaining linguistic patterns; accountability to disciplinary standards of evidence and argument; and results intended to bear on linguistic capacities, practices, systems, histories, or products.

The small distinction matters. Claims about \enquote{linguists} may show what follows for Reiss's Chomskyan subkind. They don't automatically show what follows for \term{linguists}. The bridge-principle demand applies reflexively: the argument has to explain why conclusions licensed inside an I-language programme project onto corpus linguists, typologists, historical linguists, sociolinguists, language philosophers, descriptive fieldworkers, or construction grammarians.

Without that bridge, the argument doesn't establish that non-UG \term{linguists} lack grounds for comparison. It establishes only that they don't use Reiss's preferred licence.

The boundary-policing language can then be treated methodologically rather than personally. The objection isn't that Reiss uses a stipulation. The problem is that the paper sometimes lets a local stipulation do global work.

\section{Projectibility as an Alternative Licence}

A projectibility account starts with the projection target. What does the proposed category let us infer? Which property-and-relation pattern supports that inference? Which diagnostics identify the pattern? What ordering relation makes source and target properties co-available rather than accidentally associated? What stabilizes that co-availability? Which scope conditions and failure modes mark the projection's limits?

In this vocabulary, a \term{profile} is the worldly property-and-relation pattern that makes a projection worth trying. A \term{profile specification} is the inquiry-facing account that says how the profile is located, what support grade it has, what it projects to, and where it fails \citep{reynolds2026kindsProjectibilityProfiles}. The distinction matters because a coding rule, a corpus cue, or a judgement contrast may locate a profile without explaining why the projected property follows.

Support comes in grades. A diagnostic cue identifies a candidate contrast but doesn't by itself license transfer. A scoped bridge specifies a target, profile, diagnostics, ordering relation, scope, and defeaters; it licenses a local projection. A robust bridge earns broader use only when it survives independent cases, held-out contrasts, and known defeaters. These grades aren't decorations. They're how a bridge principle avoids becoming an all-purpose permission slip.

A simple register example shows the scale. A corpus frequency difference may cue a register contrast. It becomes a scoped bridge only when the target is specified, the conditioning structure is identified, and the projection survives held-out material: for instance, knowing the institutional setting predicts likely forms, and hearing the forms helps infer the setting. It becomes robust only if the same profile continues to support projections across genres, speakers, institutions, or periods while respecting declared defeaters. The point isn't to inflate every coding decision into metaphysics. The point is to keep cue, warrant, and scope from collapsing into one another.

Projectibility is a second-order discipline for evaluating proposed bridges, not a rival language faculty. UG is one possible first-order filling of the schema: for some I-language target, the proposed profile may be ordered and stabilized by a genetically constrained language faculty. The claim here is that UG doesn't get monopoly status by being the only familiar filling. Functional-typological recurrence, diachronic pathway, communicative recurrence, processing pressure, acquisition bottleneck, social uptake, and institutional maintenance can also fill the schema when they explain co-availability for the declared target.

Haspelmath's comparative concepts already supply one non-UG discipline for crosslinguistic work. A comparative concept is constructed for crosslinguistic comparison rather than discovered as a universal category instantiated in each language \citep{haspelmath2010}. That view doesn't make comparison impossible. It makes comparison accountable: the comparative concept has to be defined, its diagnostics have to be stated, and language-particular categories can't simply be equated because they share a familiar label.

Croft's Radical Construction Grammar supplies another. Categories are language-specific, but constructions can still be compared through function and typological patterning \citep{croft2001}. The bridge Reiss's monopoly claim rules out too quickly is comparison without a universal inventory of Chomskyan primitives.

Comparanda/realization discipline sharpens the point \citep{reynolds2025comparanda}. A comparative concept is a definition and diagnostic procedure for a study. A comparandum is the target pattern placed under comparison. Language-internal realizations are the forms, constructions, functions, or categories that carry that pattern in a particular grammar. Cross-level couplings among those pieces are bridge hypotheses, supported or rejected by diagnostics and predictions, not identities between levels.

Three operations have to be kept apart. Comparison places phenomena under a declared description. Classification assigns a language-internal realization to a comparandum. Evidential transfer treats evidence about one realization as support for an analysis of another. Comparative concepts often license the first operation and discipline the second. They don't by themselves license the third. Evidential travel begins only when the comparative setup is embedded in a profile with projection, support, scope, and defeaters.

Haspelmath's instrumental caution is compatible with this view. Some comparative concepts are constructed tools. A profile claim is stronger only when the tool tracks a worldly recurrence that supports projections for a purpose. That recurrence needn't be a genetically specified UG primitive. It may be functional, distributional, diachronic, normative, or institutional, and it remains defeasible. The account can be realist without treating every useful comparandum as an innate category.

Bridge-principle terms make the same point. A bridge principle is a scoped licence from observed diagnostics to a projected target through a proposed profile: if source features S identify profile P under conditions C, then target expectation Q is licensed until defeaters D appear. The licence depends on the target, and the target fixes the relevant explanatory level.

For linguistic work, order and stabilization can come apart. Formal dependency may order a grammar-internal projection. Acquisition pressure, processing bias, diachronic pathway, communicative recurrence, social uptake, or institutional maintenance may stabilize a profile across a declared range. None of those is a magic pass. Each has to explain why observed source properties and projected target properties are co-available in that case.

These alternatives aren't interchangeable. Some are weak, local, or task-specific. That's an advantage. A serious bridge principle should be able to demote itself: localize to a target, lower its support grade, or reject transfer when the expected projection fails. UG can remain a candidate bridge where the projection target is an individual mental grammar and the proposed stabilizing support is a genetically constrained language faculty. It doesn't follow that UG is the only bridge for every linguistic target.

\section{A Worked Audit: French Evidence for English Polarity}

Reiss's own French-English example is the best test case. English \mention{any} occurs in negative-polarity environments and in free-choice uses. French, in the relevant standard examples, uses forms such as \olang{rien} in negative contexts and \olang{n'importe quoi} in free-choice contexts. Reiss uses the French contrast to argue that English has two forms pronounced \mention{any}, and then asks what licenses French evidence for English. His answer is UG \citep[sec.~5.3.3]{reiss2026armchair}.

Projectibility-first reconstruction keeps the useful inference while denying the monopoly. The target isn't \enquote{English as such}; it's the analysis of English \mention{any}: whether a single surface form realizes one lexical-semantic item or two. The comparandum is the functional contrast between polarity-dependent uses and free-choice uses. The language-internal realizations include English \mention{any}, French \olang{rien}, and French \olang{n'importe quoi}. Classification of English \mention{any} as one item or two is exactly what has to be argued for; it can't be read off the French split.

Diagnostics remain local. Downward-entailing contexts, negation, questions, conditionals, and related environments diagnose the negative-polarity side. Free-choice readings, modal environments, and indifference readings diagnose the other side. Translation into French is a diagnostic cue because it shows that the semantic-pragmatic contrast can be lexicalized separately, not because it transfers a French category into English.

A simpler case shows why that restriction matters. French expresses age with \olang{avoir}: \olang{j'ai cinquante-sept ans}. English has \mention{have}, numerals, and \mention{years}, and \mention{I have five years} is grammatical when it describes a duration, as in time remaining until retirement. But \mention{I have fifty-seven years} isn't a native English age predication. French evidence can bear on English here only after the target is named: age predication, duration expression, calque, bilingual transfer, or translation practice. The French construction is a cue to a contrast, not a licence to import the French analysis into English.

That caution isn't pedantry. Polarity-item classifications are famously not exhausted by a single label or a simple binary contrast. Hoeksema's survey of negative polarity items shows many distributional classes across questions, comparatives, conditionals, universals, \emph{only}, superlatives, and related environments: the patterns aren't random, but neither do they reduce to one strong/weak taxonomy \citep{hoeksema2012}. Reiss's example needs exactly the sort of item-sensitive distributional profile that a projectibility audit demands. The label \enquote{NPI} doesn't do the transfer work.

French creates a second problem. French \olang{rien} is entangled with negative-quantifier and n-word behaviour, not simply with the behaviour of English NPI \mention{any} as treated in classic polarity accounts \citep{ladusaw1979,chierchia_logic_2013}. If \olang{rien} is made to stand for the polarity-dependent member of an English split, then a classification has already been imported: negative quantification, negative concord, polarity sensitivity, and free choice have been lined up in the way the argument requires. That's not a bridge principle. It's the very classificatory shortcut Reiss says his opponents aren't entitled to make. The bridge has been built by assuming the destination classification it's supposed to license.

Mere parsimony can't be the ordering relation. The ordering relation has to explain why the observed distributions and the projected analysis are co-available. In this case, the candidate relation is compositional and distributional: polarity-dependent readings are tied to licensing environments, while free-choice readings carry modal or indifference force in environments where that licensing relation is absent. If English-internal evidence shows that those two dependency profiles behave differently across environments, a split analysis gains support. If it doesn't, French remains a prompt rather than a licence.

Nor is the stabilizer the evidence itself. Overt splits in other languages and English distributional contrasts are diagnostics. Candidate stabilizers would have to be defended separately: recurrent functional pressure to keep polarity-dependent indefinites apart from free-choice expressions, a diachronic pathway that repeatedly separates the two, a processing or acquisition pressure, or a grammatical theory that predicts the split. Without such support, the French evidence has cue status. With English-internal diagnostics and a defended recurrence profile, it may become a scoped bridge. It isn't a robust bridge on Reiss's example alone.

Defeaters matter just as much. The bridge weakens if English-internal diagnostics fail to separate the uses, if a unified semantics predicts both distributions without residue, if the French contrast turns out to reflect idiosyncratic lexicalization rather than the relevant profile, or if dialect and register variation in French undermine the claimed contrast. Reiss himself notes complications involving negative concord, French dialect variation, and the mismatch between standard written French and spoken varieties. A projectibility-first account treats those complications as defeaters or scope conditions, not as footnotes to a universal permission to transfer.

That reconstruction gives a verdict different from the bare UG postulate. Reiss treats the French data as licensed for English by UG. Projectibility demotes the French data to a cue until the target, diagnostics, ordering relation, stabilizer, scope, and defeaters are supplied. A substantive UG theory may supply some of that missing structure, but then it's functioning as one first-order bridge under audit. The general postulate alone hasn't licensed the transfer. Reiss's showcase case needs the very bridge machinery his monopoly conclusion tries to reserve for UG.

\section{A Successful Non-UG Bridge: Historical Comparison}

The audit of French-English \mention{any} is deliberately severe: it shows that Reiss's own example doesn't reach robust-bridge status on the evidence supplied. But non-UG bridges needn't stop at cue level. Historical-comparative linguistics provides a familiar case in which evidence from one language bears on another because the bridge is causal-historical rather than UG-theoretic.

Here, the target is an inherited pattern: a form, sound correspondence, morphological alternation, construction, or subgroup relation in a family of languages. The profile is common descent under transmission and change. Diagnostics include recurring cognate sets, systematic sound correspondences, shared irregular morphology, chronology, geography, and independent evidence about borrowing and contact \citep{fortson2004}. The ordering relation is causal: daughter systems inherit, transform, lose, borrow, analogize, and reanalyse material from earlier systems. Evidence from Sanskrit, Greek, Latin, Gothic, or Old Irish can bear on a reconstruction not because UG says all languages are comparable, but because those languages are historically related in ways that constrain what counts as a good projection.

A miniature case makes the bridge visible. A correspondence set such as Latin \olang{pater}, Greek \olang{pat\={e}r}, Sanskrit \olang{pitar-}, and Gothic \olang{fadar} doesn't merely record resemblance. Together with many comparable sets, it supports a reconstruction of an inherited kinship term and a correspondence profile in which Germanic \olang{f} aligns with non-Germanic \olang{p} under independently motivated sound change \citep{fortson2004}. Evidence from Gothic can bear on the reconstruction, and the reconstruction can bear on the analysis of the Latin, Greek, and Sanskrit forms, because transfer is mediated by a causal-historical profile.

This bridge has real scope limits. It licenses evidential travel within a family or contact history; it doesn't license arbitrary transfer from any language to any other. Its defeaters are also familiar: chance resemblance, borrowing, areal diffusion, analogy, expressive vocabulary, semantic drift, defective attestation, and chronology that makes inheritance impossible. Those defeaters aren't embarrassments. They're what make the bridge scientific. A claim about inherited structure can be localized, demoted, or rejected when the correspondence profile fails.

UG may still matter to historical linguistics in other ways: possible sound systems, learnability, or constraints on grammar change may set background conditions. But that isn't the bridge that lets evidence from one Indo-European language bear on another. The transfer is licensed by a reconstruction profile with declared diagnostics, ordering, scope, and defeaters. That's exactly what Reiss's monopoly claim says non-UG \term{linguists} lack, and it's exactly what a major linguistic tradition has long used. The comparative method licensed this sort of evidential travel long before UG was a Chomskyan postulate. Either that tradition already had a non-UG bridge, or a century of successful reconstruction wasn't linguistic evidence at all.

Historical comparison kills the monopoly claim, but it also marks a limit. It's a vertical bridge: evidence travels because languages stand in a causal-genealogical relation. Reiss's motivating examples include horizontal transfer: Japanese evidence on English, or typological evidence across unrelated languages, where common descent is unavailable by construction. The historical case settles one issue and leaves another open. Non-UG bridge principles exist, so UG can't be the exclusive licence for linguistic evidential travel.

For unrelated-language transfer specifically, the French \mention{any} audit supplies a method and a demotion verdict; Haspelmath and Croft discipline comparison and classification; a successful horizontal non-UG bridge would need its own profile, stabilizers, held-out projections, and defeaters. The claim here is the negative one Reiss denies: UG has no monopoly. The stronger positive typological case is a further project rather than a hidden premise.

\section{Public and Distributed Linguistic Objects}

So far, the paper has answered the evidential half of the monopoly claim. Reiss also presses an individuation worry: public languages and standards aren't compact mental objects, and Standard English doesn't travel to the moon with Neil Armstrong. That worry is real. The reply shouldn't be that public languages are I-languages after all, or that individuation doesn't matter. The reply is that public, community-level, and normatively regimented objects can be individuated by the projections they support and the failure conditions they accept. Their tractability can come from maintenance, uptake, correction, transmission, institutionalization, and projectibility.

Stainton's defence of public languages makes this point without denying mental grammars \citep{stainton2011public}. The study of public languages can be viable if the object is understood as a humanly individuated system whose evidence base includes deployment, uptake, and psychological facts. Santana's pluralist treatment of language also cuts against the idea that one privileged language-object settles all legitimate inquiry \citep{santana2016language}.

Reiss's remarks about Standard English and public languages depend on a serious individuation worry. But the moon example applies a luggage test to a distributed social object. A public standard isn't a medium-sized object that has to travel to the moon with Neil Armstrong, but neither does it become an object of inquiry by stipulation. It has to be individuated by the role it plays in projections. Model-theoretic syntax makes related room for formally characterizing public constructs without treating them as mental grammars \citep{pullum2001}; sociolinguistic work on standardization and enregisterment supplies further mechanisms for norm uptake and public recognizability \citep{milroy1999,agha2005}. The point isn't that Standard English and I-language are rival objects of the same kind. The point is that formal, public, and psychological targets answer different questions.

Once the explanatory level is named, the apparent contrast loses its force. I-language inquiry asks one set of questions about individual mental grammars. Public-language inquiry asks another set about uptake, correction, transmission, standardization, and shared norms. Corpus inquiry asks another set about distributions in attested records. Typology asks another set about cross-system recurrence. Calling only the first target \enquote{language} doesn't make the other targets disappear.

This isn't another example of crosslinguistic evidential transfer. It answers the ontological half of Reiss's restriction. For non-I-language targets, the bridge-principle demand still applies, but the relevant ordering relation and stabilizers may be social-practical rather than genetic.

Varieties give a compact test case. Register, dialect, and discourse community can be analysed as conditioning structures: participant expectation conditioned primarily on situation, ascription, or identification \citep{reynolds2026varieties}. The bridge claim is projective. Knowing enough about the conditioning state should help predict likely forms, and hearing the forms should help infer the conditioning state. When those predictions fail, the variety claim localizes or demotes rather than disappearing into metaphysical fog.

Standard written English fits the same template. Its target is an edited public standard, not every English speaker's mental grammar or every English utterance. The projected expectations concern spelling, punctuation, selected morphosyntactic forms, correction practices, uptake in institutional settings, and the treatment of deviations as errors, choices, or markers of another register. The diagnostics include style guides, editorial intervention, school assessment, institutional documents, dictionaries, usage manuals, and corpora of edited prose.

Here, the ordering relation is normative and institutional rather than genetic: correction, publication, assessment, and uptake partly make the standard whose future behaviour they help predict. That constitutive relation isn't a hidden mechanism; it's the relevant ordering relation for a social-practical target.

For example, evidence from style guides, corrected student prose, copy-edited publications, and edited corpora can support a projection about how the \mention{its}/\mention{it's} distinction will be treated in academic prose. It doesn't support the same projection about spontaneous speech, a speaker's I-language, or every digital text. The target controls the inference.

The same template applies beyond orthography. Institutional treatment of \mention{between you and I}, \mention{who}/\mention{whom}, preposition stranding, or singular \mention{they} can support projections about classification, correction, acceptance, or style-shifting in edited prose. The projection isn't what every speaker's I-language contains. It's how a form will be treated in an institutionally answerable public standard.

Schooling, publishing, examination, professional editing, reference works, and repeated public use stabilize the profile. The scope is edited and institutionally answerable prose. It isn't casual speech, all written English, all Englishes, or the competence of every individual speaker. Defeaters include divergent institutional norms, stable acceptance of a formerly corrected form, failure to predict held-out edited texts, or evidence that a supposed standard feature tracks genre, medium, or identity rather than the standard.

The template licenses some projections and blocks others. It explains why evidence from a style guide may bear on an edited newspaper sentence but not on a child's spontaneous speech; why corpus frequencies in unedited social media may be diagnostic for one target and misleading for another; and why public-language inquiry doesn't need to pretend that public standards are located in one head. A distributed standard is disciplined when its target, diagnostics, stabilizers, scope, and defeaters are declared.

Those are bridge claims, not slogans. Uptake, correction, transmission, and institutional maintenance may be stabilizing backgrounds, parts of the profile, or mere diagnostics depending on the projection target. A public-language argument has to say which role they play. That demand is a strength of the alternative, because it imposes the same discipline on non-UG objects that Reiss rightly imposes on crosslinguistic inference.

\section{Conclusion}

Crosslinguistic inference requires bridge principles; it doesn't require Chomskyan Universal Grammar. Reiss's paper is strongest when it insists that categories are theory-mediated and that evidence doesn't travel for free. It's weakest where it treats that methodological burden as if it had only one possible answer.

Historical reconstruction also shows why the projectibility claim is realist rather than merely instrumental. Reconstructed proto-forms aren't convenience labels. They're posits constrained by converging correspondences, causal order, independent chronology, and defeaters. Public standards and typological comparanda are less compact objects, but they answer to the same demand: declare the projections they support and the failures they accept.

Projectibility-first work accepts the burden and raises it. For every proposed evidential transfer, say what the target is, which profile carries the inference, which diagnostics identify it, what orders and stabilizes it, where the inference stops, and what would defeat it. Reiss's \enquote{linguists} can do that for a Chomskyan I-language programme. Historical-comparative \term{linguists} can do it with common descent and reconstruction. Public-language researchers can do it with uptake, correction, and institutional maintenance. Reiss hasn't shown that \term{linguists} need UG to do linguistics.

Failure conditions also apply to projectibility. The account would fail if non-UG bridges repeatedly collapsed into cue-level prompts, if proposed profiles couldn't predict held-out cases, if ordering relations were indistinguishable from labels for the data, or if social and typological targets produced no projections beyond the diagnostics used to identify them. Those are real risks. They're also the risks Reiss's UG monopoly avoids only by not asking when its own bridge should demote.

\section*{Acknowledgements}

OpenAI Codex and Anthropic Claude Opus served as drafting, review, and editing aids during the preparation of this paper. I am responsible for all theoretical claims, arguments, errors, and interpretive choices.

\clearpage
\printbibliography

\end{document}
